%%=============================================================================
%% Samenvatting
%%=============================================================================

% TODO: De "abstract" of samenvatting is een kernachtige (~ 1 blz. voor een
% thesis) synthese van het document.
%
% Een goede abstract biedt een kernachtig antwoord op volgende vragen:
%
% 1. Waarover gaat de bachelorproef?
% 2. Waarom heb je er over geschreven?
% 3. Hoe heb je het onderzoek uitgevoerd?
% 4. Wat waren de resultaten? Wat blijkt uit je onderzoek?
% 5. Wat betekenen je resultaten? Wat is de relevantie voor het werkveld?
%
% Daarom bestaat een abstract uit volgende componenten:
%
% - inleiding + kaderen thema
% - probleemstelling
% - (centrale) onderzoeksvraag
% - onderzoeksdoelstelling
% - methodologie
% - resultaten (beperk tot de belangrijkste, relevant voor de onderzoeksvraag)
% - conclusies, aanbevelingen, beperkingen
%
% LET OP! Een samenvatting is GEEN voorwoord!

%%---------- Nederlandse samenvatting -----------------------------------------
%
% TODO: Als je je bachelorproef in het Engels schrijft, moet je eerst een
% Nederlandse samenvatting invoegen. Haal daarvoor onderstaande code uit
% commentaar.
% Wie zijn bachelorproef in het Nederlands schrijft, kan dit negeren, de inhoud
% wordt niet in het document ingevoegd.

\IfLanguageName{english}{%
    \selectlanguage{dutch}
    \chapter*{Samenvatting}
    \chapter{\IfLanguageName{dutch}{Samenvatting}{Abstract}}
    
    Security Operations Centers (SOC’s) worden dagelijks geconfronteerd met een grote hoeveelheid security-events en alerts afkomstig van uiteenlopende databronnen. Traditionele Security Information and Event Management (SIEM)-systemen steunen voornamelijk op rule-based correlatie, wat beperkingen vertoont bij complexe en meerlagige aanvalspatronen. Hierdoor ontstaat een hoge werkdruk voor analisten en een verhoogd risico op gemiste dreigingen.
    \newline
    \newline
    Deze bachelorproef onderzoekt hoe graph-based machine learning (GBML) kan worden toegepast om security-events automatisch te correleren en te chainen. Door events en entiteiten te modelleren als nodes en hun onderlinge relaties als edges binnen een graafstructuur, wordt beoogd om contextuele verbanden expliciet te maken en multi-stage aanvalspatronen beter te detecteren.
    \newline
    \newline
    Op basis van een literatuurstudie worden de beperkingen van traditionele correlatiemechanismen geanalyseerd en worden graph-gebaseerde technieken onderzocht als mogelijke oplossing. Vervolgens werd een methodologie uitgewerkt en geïmplementeerd in de vorm van een proof-of-concept, waarbij security-events worden omgezet naar een grafenstructuur en geanalyseerd met een Relational Graph Convolutional Network (R-GCN).
    \newline
    \newline
    De aanpak werd geëvalueerd op twee datasets met een verschillende doelstelling. Voor de eerste dataset werd een supervised multiclass classificatiemodel opgebouwd. Het R-GCN model behaalde hierbij een accuracy van 0.8738 en een weighted F1-score van 0.8744, wat aantoont dat de graph-based representatie bruikbaar is voor het onderscheiden van verschillende aanvalstypes.
    \newline
    \newline
    Voor de tweede dataset werd gewerkt in een weakly supervised setting waarbij de focus verschoof van pure classificatie naar prioritisatie, ranking en chaining van verdachte events. Door de sterke class imbalance en beperkte labels bleken klassieke classificatiemetrics minder representatief, waardoor bijkomende aandacht werd besteed aan de interpretatie en contextualisatie van events.
    \newline
    \newline
    De resultaten tonen dat graph-based machine learning zowel inzetbaar is voor classificatie als voor contextuele analyse van security-events.. Afhankelijk van de beschikbare data verschuift de rol van het model van detectie naar ondersteuning van prioritisatie en chaining.
    \selectlanguage{english}
}{}

%%---------- Samenvatting -----------------------------------------------------

\chapter*{\IfLanguageName{dutch}{Samenvatting}{Abstract}}

Security Operations Centers (SOC’s) worden dagelijks geconfronteerd met een grote hoeveelheid security-events en alerts afkomstig van uiteenlopende databronnen. Traditionele Security Information and Event Management (SIEM)-systemen steunen voornamelijk op rule-based correlatie, wat beperkingen vertoont bij complexe en meerlagige aanvalspatronen. Hierdoor ontstaat een hoge werkdruk voor analisten en een verhoogd risico op gemiste dreigingen.
\newline
\newline
Deze bachelorproef onderzoekt hoe graph-based machine learning (GBML) kan worden toegepast om security-events automatisch te correleren en te chainen. Door events en entiteiten te modelleren als nodes en hun onderlinge relaties als edges binnen een graafstructuur, wordt beoogd om contextuele verbanden expliciet te maken en multi-stage aanvalspatronen beter te detecteren.
\newline
\newline
Op basis van een literatuurstudie worden de beperkingen van traditionele correlatiemechanismen geanalyseerd en worden graph-gebaseerde technieken onderzocht als mogelijke oplossing. Vervolgens werd een methodologie uitgewerkt en geïmplementeerd in de vorm van een proof-of-concept, waarbij security-events worden omgezet naar een grafenstructuur en geanalyseerd met een Relational Graph Convolutional Network (R-GCN).
\newline
\newline
De aanpak werd geëvalueerd op twee datasets met een verschillende doelstelling. Voor de eerste dataset werd een supervised multiclass classificatiemodel opgebouwd. Het R-GCN model behaalde hierbij een accuracy van 0.8738 en een weighted F1-score van 0.8744, wat aantoont dat de graph-based representatie bruikbaar is voor het onderscheiden van verschillende aanvalstypes.
\newline
\newline
Voor de tweede dataset werd gewerkt in een weakly supervised setting waarbij de focus verschoof van pure classificatie naar prioritisatie, ranking en chaining van verdachte events. Door de sterke class imbalance en beperkte labels bleken klassieke classificatiemetrics minder representatief, waardoor bijkomende aandacht werd besteed aan de interpretatie en contextualisatie van events.
\newline
\newline
De resultaten tonen dat graph-based machine learning zowel inzetbaar is voor classificatie als voor contextuele analyse van security-events. Afhankelijk van de beschikbare data verschuift de rol van het model van detectie naar ondersteuning van prioritisatie en chaining.